\fiche{1}{Prédire le suivant : la chaîne complète}{\src{05_model.py}{MiniGPT.forward}}
\objectif{Comprendre ce que calcule un modèle de langage avant d'étudier ses opérations.}

Un texte arrive sous forme de caractères, non de nombres réels. Un
\emph{tokeniseur} le transforme en une suite d'identifiants entiers. Dans le
mode caractère du dépôt, un token représente généralement un caractère ;
dans le mode BPE, il peut représenter un fragment différent. Le vocabulaire
est l'ensemble des tokens disponibles et sa taille se note \(V\). Un
identifiant est une étiquette : le token numéro 8 n'est pas deux fois
plus important que le token numéro 4.

La tâche est de prédire le token immédiatement suivant à partir des tokens
déjà visibles. Pour une suite fictive \((2,5,1,4)\), une fenêtre fournit
l'entrée \((2,5,1)\) et les cibles \((5,1,4)\). À la première position,
le modèle doit proposer 5 en ne connaissant que 2 ; à la deuxième,
il dispose du préfixe \((2,5)\). Le mécanisme causal empêche de consulter
les positions futures, même lorsqu'elles sont présentes dans le tableau.

\notion{La chaîne de calcul.}
Dans \code{MiniGPT.forward}, \code{token\_embedding} transforme chaque
identifiant en une liste de \(D\) nombres, appelée représentation.
Le mode positionnel appris ajoute \code{position\_embedding}. Puis
\code{blocks}, \code{final\_norm} et \code{lm\_head} produisent des scores,
les \emph{logits}. Le mode RoPE traite autrement les positions ; cela
ne change pas l'objectif prédictif.
\[
 \text{texte}\longrightarrow\text{identifiants}
 \longrightarrow\text{représentations}
 \longrightarrow\text{logits}\longrightarrow\text{probabilités}.
\]
Chaque position produit \(V\) scores, pas un identifiant directement.
Une transformation appelée softmax donnera des probabilités positives
dont la somme vaut un. À l'entraînement, on mesure la probabilité attribuée
à la cible connue. En génération, on choisit un token, on l'ajoute à la
suite, puis on recommence. Ces deux usages partagent le même prédicteur.

Nous noterons \(B\) le nombre de fenêtres traitées ensemble, \(T\) leur
longueur courante et \(C\) la capacité maximale de contexte. Sans cache,
on doit avoir \(1\leq T\leq C\). Les \(L\) blocs transforment les
représentations sans changer le nombre de positions. Aucun de ces nombres
ne désigne une quantité de mots comprise par la machine.

\exercice{Avec \(B=2\), \(T=3\) et \(V=4\), combien de scores et combien
de cibles un passage produit-il ? Quelle cible correspond au dernier
token d'entrée de notre exemple ?}
\corrige{Chaque fenêtre fournit trois listes de quatre scores, donc
\(2\times3\times4=24\) scores pour six cibles. La dernière entrée vaut 1,
mais sa cible vaut 4 : recopier l'entrée ne résoudrait pas la tâche.
Les probabilités portent sur quatre possibilités à chacune des six
positions, non sur les six positions elles-mêmes.}

\fiche{2}{Fonctions, arguments et indices Python}{\src{02_dataset.py}{NextTokenDataset}}
\objectif{Traduire les notations mathématiques en accès à une suite Python sans décalage accidentel.}

Une fonction associe une sortie à une entrée autorisée. Au lycée,
\(f(x)=2x+1\) signifie qu'une même recette s'applique à chaque valeur
de \(x\). Un programme peut recevoir une liste plutôt qu'un seul nombre.
\code{SimpleTokenizer.encode}, la méthode d'encodage du tokeniseur,
reçoit un texte et renvoie
une liste d'entiers. Son domaine n'est donc pas celui d'une fonction
polynomiale, mais l'idée d'une transformation bien définie reste identique.

Une suite mathématique peut commencer à l'indice zéro : écrivons
\(s_0,s_1,\ldots,s_{N-1}\) pour \(N\) tokens. Python adopte précisément
cette convention. \code{s[0]} est le premier élément, tandis que
\code{s[N-1]} est le dernier. Les crochets indiquent un accès, pas une
multiplication. Un indice compte les positions ; la valeur rangée à cette
position est une information distincte.

\notion{Les tranches excluent leur borne de droite.}
L'expression \code{s[a:b]} prend les indices \(a,a+1,\ldots,b-1\).
Si les bornes restent dans la liste et si \(b\geq a\), elle contient
\(b-a\) éléments. Dans \code{NextTokenDataset}, la longueur de fenêtre
est \code{context\_length}, que nous notons \(C\).
\[
 x_t=s_{i+t},\qquad y_t=s_{i+t+1},
 \qquad 0\leq t<C.
\]
Pour \(s=(4,7,2,9,3)\), \(i=1\) et \(C=3\), l'entrée utilise les
indices 1, 2, 3 : elle vaut \((7,2,9)\). La cible utilise les indices
2, 3, 4 : elle vaut \((2,9,3)\). Les deux listes ont la même longueur,
mais leurs contenus sont décalés.

Lorsqu'on écrit une formule, annoncer le début des indices évite
une ambiguïté fréquente : le « troisième élément » est ici
celui d'indice deux, pas celui d'indice trois.

Il faut un token supplémentaire pour connaître la dernière cible.
Avec un pas de déplacement égal à un, les débuts admissibles vont
de zéro à \(N-C-1\), inclus : cela fait \(N-C\) fenêtres.
Ne confondons pas cette borne maximale avec le nombre de fenêtres.
Le jeu de données sur disque ajoute un paramètre de pas ; son comptage
tient aussi compte de ce paramètre.

Enfin, \code{logits[:, -1, :]} emploie un indice négatif pour sélectionner
la dernière position. Les deux deux-points conservent les autres axes :
ce n'est pas une sélection du dernier token du vocabulaire.

\exercice{Pour \(s=(8,1,6,2,5,7)\), \(C=2\) et \(i=2\), donner
\code{s[i:i+C]}, la cible et le nombre total de fenêtres.}
\corrige{L'entrée vaut \((6,2)\), la cible \((2,5)\).
Les indices de début possibles sont \(0,1,2,3\), soit \(6-2=4\)
fenêtres. Le début 4 serait impossible : sa deuxième cible demanderait
l'élément d'indice 6, absent de cette suite.}

\fiche{3}{Sommes, moyennes et pondérations}{\src{06_train.py}{evaluate}}
\objectif{Lire le symbole somme et comprendre pourquoi une moyenne de moyennes peut tromper.}

Le signe \(\sum\), lettre grecque sigma, abrège une addition répétée.
L'indice écrit en dessous parcourt les valeurs indiquées ; chaque terme
est calculé avant d'être ajouté. Par exemple,
\(\sum_{i=0}^{3}a_i=a_0+a_1+a_2+a_3\).
Cet indice est local : remplacer partout \(i\) par \(j\) ne modifie
pas le résultat. En revanche, changer les bornes change les termes inclus.

La moyenne arithmétique de \(n>0\) valeurs partage leur somme en
\(n\) parts égales. Pour les erreurs \(1,2,3,6\), la somme vaut 12
et la moyenne vaut 3. Elle a la même unité que les valeurs initiales.
La somme augmente généralement quand on ajoute des observations ;
la moyenne permet de comparer des ensembles de tailles différentes.
\[
 \bar a=\frac{1}{n}\sum_{i=0}^{n-1}a_i,
 \qquad
 \sum_{i=0}^{n-1}(a_i-\bar a)=0.
\]
La deuxième identité se vérifie en développant :
la somme des \(a_i\) vaut \(n\bar a\), que l'on soustrait à elle-même.
Elle ne signifie pas que les écarts individuels sont tous nuls.

\notion{Compter les observations plutôt que les lots.}
Dans \code{evaluate}, \code{loss.item()} représente, avec le critère
utilisé par l'entraînement, une erreur moyenne sur les tokens du lot.
\code{y.numel()} compte ces tokens. Le programme multiplie les deux,
additionne les résultats, puis divise par \code{total\_tokens}.
Un lot final plus petit reçoit ainsi un poids proportionnel à sa taille.

Supposons un premier lot de six cibles avec une erreur moyenne de 2,
puis un lot de deux cibles avec une moyenne de 6. Les erreurs totales
valent respectivement 12 et 12. La moyenne globale est donc
\((12+12)/8=3\), et non \((2+6)/2=4\).
Cette dernière opération donnerait autant d'importance à deux cibles
qu'à six.

Plus généralement, une moyenne pondérée utilise des coefficients positifs
dont la somme vaut un. Ici, les poids sont \(6/8\) et \(2/8\).
Nous retrouverons cette idée avec les probabilités et, plus tard,
avec les mélanges de représentations. Les poids n'ont pas besoin
d'être égaux pour former une moyenne.

\exercice{Trois lots contiennent respectivement 4, 4 et 2 tokens.
Leurs pertes moyennes valent 1, 3 et 5. Calculer la perte globale.}
\corrige{La somme des pertes vaut \(4\times1+4\times3+2\times5=26\).
Le nombre de tokens vaut 10, donc la moyenne correcte est \(2{,}6\).
La moyenne non pondérée des trois valeurs serait 3 : elle surestime
ici l'importance du dernier petit lot.}

\fiche{4}{Vecteurs : des listes de nombres organisées}{\src{05_model.py}{MiniGPT.forward}}
\objectif{Manipuler une représentation de token sans supposer de connaissances d'algèbre linéaire.}

Un vecteur de dimension \(D\) est ici une liste ordonnée de \(D\)
nombres réels. Dire \(u\in\mathbb{R}^{D}\) signifie simplement
\(u=(u_0,\ldots,u_{D-1})\), avec une valeur réelle à chaque place.
L'ordre compte : \((1,2)\) et \((2,1)\) ne sont pas le même vecteur.
Le mot dimension compte les coordonnées, non les tokens ou les couches.

Dans \code{MiniGPT}, \code{token\_embedding} contient une représentation
apprise pour chaque identifiant. Si \(D=3\), un token pourrait recevoir
\((0{,}2,-0{,}1,0{,}7)\). Ces coordonnées ne sont pas imposées comme
« sujet », « verbe » et « adjectif ». Elles sont ajustées ensemble
par l'apprentissage ; leur interprétation ne se lit pas directement
dans leurs indices.

\notion{Deux opérations élémentaires.}
Pour additionner deux vecteurs de même dimension, on additionne leurs
coordonnées correspondantes. Pour multiplier un vecteur par un nombre,
on multiplie toutes ses coordonnées par ce nombre, appelé scalaire.
\[
 (1,-2,3)+(4,5,-1)=(5,3,2),
 \qquad 2(1,-2,3)=(2,-4,6).
\]
Une addition ne concatène donc pas les listes : sa dimension reste 3.
Au contraire, les placer bout à bout donnerait six coordonnées et
constituerait une autre opération.

Soustraire deux vecteurs fonctionne de la même manière.
Le vecteur différence décrit un déplacement coordonnée par coordonnée ;
sa norme mesure la distance entre les deux représentations.

Dans le mode positionnel appris, \code{tok\_emb + pos\_emb} ajoute
un vecteur lié au token et un vecteur lié à sa position.
Pour un token représenté par \((1,0,-1)\) et une position représentée
par \((0{,}1,0{,}2,0{,}3)\), le résultat vaut
\((1{,}1,0{,}2,-0{,}7)\). Le même token placé ailleurs garde sa
représentation initiale, mais reçoit une autre contribution positionnelle.

La longueur géométrique, ou norme euclidienne, généralise Pythagore :
\(\|u\|=\sqrt{u_0^2+\cdots+u_{D-1}^2}\).
Ainsi \(\|(3,4)\|=5\). Cette longueur mesure une amplitude,
pas une probabilité. Des coordonnées négatives et des normes supérieures
à un sont parfaitement autorisées pour les représentations du modèle.

Il est utile de séparer trois notions : l'identifiant entier choisit
une ligne de table ; le vecteur représente ce token ; sa norme résume
une seule propriété de ce vecteur. Deux vecteurs de même norme peuvent
contenir des informations très différentes.

\exercice{Un token a pour représentation \(u=(2,-1)\), sa position
\(p=(-1,3)\). Calculer \(h=u+p\), \(h/2\), puis \(\|h\|\).}
\corrige{On obtient \(h=(1,2)\), puis \(h/2=(0{,}5,1)\).
Sa norme vaut \(\sqrt{1^2+2^2}=\sqrt5\).
Le vecteur \(h\) possède toujours deux coordonnées : l'addition
positionnelle n'a ni ajouté un token ni doublé la dimension.}

\fiche{5}{Produit scalaire : comparer et combiner}{\src{03_attention.py}{scaled_dot_product_attention}}
\objectif{Calculer un score à partir de deux vecteurs et distinguer ce score d'une probabilité.}

Le produit scalaire de deux vecteurs de même dimension consiste à
multiplier les coordonnées correspondantes, puis à additionner les
produits. Son résultat est un seul nombre, d'où son nom.
Avec \(u=(1,2,-1)\) et \(v=(3,0,4)\), les trois contributions
sont \(3,0,-4\), donc le score total vaut \(-1\).
\[
 u\cdot v=\sum_{j=0}^{D-1}u_jv_j,
 \qquad (1,2,-1)\cdot(3,0,4)=3+0-4=-1.
\]
La dimension doit être identique pour savoir quelles coordonnées
apparier. Le produit scalaire diffère du produit coordonnée par
coordonnée : ce dernier conserverait la liste \((3,0,-4)\),
alors que le produit scalaire la réduit à sa somme.

L'ordre des deux vecteurs ne change pas ce nombre, puisque chaque
produit de réels peut être inversé. En revanche, permuter les
coordonnées d'un seul vecteur modifie généralement les associations
et donc le score.

\notion{Une lecture comme combinaison pondérée.}
Si \(u\) contient des observations et \(v\) des coefficients,
le score \(u\cdot v\) mesure leur combinaison. Un coefficient nul
ignore une observation ; un coefficient négatif inverse son effet.
Le modèle apprend précisément des coefficients de ce type.
Tous les coefficients ne sont pas nécessairement positifs.

Une deuxième lecture est géométrique. Pour des vecteurs non nuls,
\(u\cdot v=\|u\|\|v\|\cos\theta\), où \(\theta\) est l'angle
entre leurs directions. Un angle droit donne zéro. Mais zéro
peut aussi venir de contributions qui s'annulent :
\((1,1)\cdot(1,-1)=1-1=0\). Il n'implique donc pas que toutes
les coordonnées soient nulles.

Les exemples du script appellent la fonction
\code{scaled\_dot\_product\_attention}. La ligne
\code{q @ k.transpose(-2, -1)} calcule de nombreux produits
scalaires à la fois. Nous étudierons leur organisation plus tard ;
pour l'instant, chaque case obtenue est simplement un score
entre deux listes de même longueur.

Un grand score n'est pas une preuve absolue de ressemblance.
Multiplier \(u\) par 10 multiplie \(u\cdot v\) par 10,
même si sa direction reste inchangée. C'est pourquoi l'amplitude
des vecteurs intervient dans ces scores. La division ultérieure
par \(\sqrt{d_k}\), avec \(d_k=D/H\) pour \(H\) têtes,
sera expliquée dans les fiches consacrées à l'attention.

\exercice{Calculer \(a\cdot b\), puis \((2a)\cdot b\), pour
\(a=(2,-1,3)\) et \(b=(1,4,0)\). Le premier résultat peut-il
être utilisé directement comme probabilité ?}
\corrige{Le premier score vaut \(2\times1+(-1)\times4+3\times0=-2\).
Le second vaut \(4-8+0=-4\), soit deux fois le premier.
Une probabilité ne peut pas être négative : ces scores nécessitent
une transformation supplémentaire avant une interprétation probabiliste.}

\fiche{6}{Matrices : lignes, colonnes et tables de paramètres}{\src{05_model.py}{MiniGPT}}
\objectif{Lire un tableau rectangulaire et relier ses dimensions à une table de représentations.}

Une matrice est un tableau rectangulaire de nombres. Sa forme indique
d'abord le nombre de lignes, puis le nombre de colonnes. Une matrice
\(A\) de forme \(2\times3\) contient donc six nombres.
Le coefficient \(A_{ij}\) se trouve à la ligne \(i\), colonne \(j\).
Nous conservons les indices commençant à zéro pour faciliter la
correspondance avec Python.
\[
 A=\begin{pmatrix}1&-2&0\\3&4&5\end{pmatrix},
 \qquad A_{0,1}=-2,\quad A_{1,2}=5.
\]
Une ligne de cette matrice est un vecteur de dimension 3.
Une colonne possède deux coordonnées. Dire seulement « une matrice
de six éléments » perd une information essentielle : les formes
\(1\times6\), \(2\times3\) et \(3\times2\) n'organisent pas
les données de la même manière.

\notion{Une matrice peut servir de table.}
Dans \code{MiniGPT}, le poids de \code{token\_embedding} a pour
forme \([V,D]\). Chaque ligne correspond à un identifiant, chaque
colonne à une coordonnée de représentation. Chercher le token
d'identifiant 2 revient à sélectionner la ligne 2, et non à
multiplier toutes les valeurs du tableau par 2.

Avec \(V=4\) et \(D=3\), cette table contient \(4\times3=12\)
paramètres. Une entrée \((2,0,2)\) consulte trois lignes, dont
deux fois la même. Les deux occurrences du token 2 reçoivent
initialement le même vecteur. Leur position et leur contexte
pourront ensuite différencier leurs représentations.

L'addition de matrices exige des formes identiques dans sa définition
élémentaire : chaque case s'ajoute à la case correspondante.
Multiplier une matrice par un scalaire multiplie toutes ses cases.
Ces opérations ne mélangent pas les colonnes ; le produit matriciel,
étudié ensuite, réalise précisément un mélange plus structuré.

Les tables ne sont pas toutes carrées. La table positionnelle apprise
a la forme \([C,D]\), alors que la table des tokens a la forme
\([V,D]\). Elles partagent la dimension des représentations mais
n'indexent pas les mêmes objets. \(C\) compte les positions
admissibles, tandis que \(V\) compte les tokens distincts.

\exercice{Une table contient les lignes \((1,0)\), \((0,2)\) et
\((-1,3)\). Donner sa forme, son nombre de paramètres et le tableau
obtenu en consultant les identifiants \((2,1)\).}
\corrige{La forme est \(3\times2\), donc la table contient six
paramètres. Les lignes sélectionnées sont \((-1,3)\), puis \((0,2)\).
Le résultat est une matrice \(2\times2\). Le nombre de lignes
de sortie dépend du nombre de tokens consultés, pas de la taille
totale du vocabulaire.}

\fiche{7}{Produit matriciel : un calcul complet}{\src{03_attention.py}{pytorch_attention_example}}
\objectif{Effectuer chaque case d'un produit et comprendre le rôle de la dimension commune.}

Multiplier deux matrices n'est pas multiplier simplement leurs cases.
Une ligne de la première rencontre une colonne de la seconde :
leur produit scalaire donne une case du résultat. Ainsi une matrice
\(A\) de forme \(m\times n\) peut multiplier une matrice \(R\)
de forme \(n\times p\). Le nombre \(n\) doit être commun ;
le résultat a la forme \(m\times p\).
\[
 A=\begin{pmatrix}1&2&-1\\0&3&2\end{pmatrix},
 \qquad R=\begin{pmatrix}2&1\\-1&0\\4&2\end{pmatrix}.
\]
Ici, deux lignes de longueur trois rencontrent deux colonnes de
longueur trois. La première case vaut
\(1\times2+2\times(-1)+(-1)\times4=-4\).
La deuxième vaut \(1\times1+2\times0+(-1)\times2=-1\).
La troisième vaut \(0\times2+3\times(-1)+2\times4=5\).
La dernière vaut \(0\times1+3\times0+2\times2=4\).
\[
 AR=\begin{pmatrix}-4&-1\\5&4\end{pmatrix},
 \qquad (AR)_{ij}=\sum_{r=0}^{n-1}A_{ir}R_{rj}.
\]
Les indices libres \(i,j\) désignent la case conservée.
L'indice \(r\), parcouru dans la somme, disparaît du résultat.
Cette différence explique pourquoi la dimension commune ne figure
plus dans la forme finale.

\notion{Le code applique cette recette en parallèle.}
L'opérateur Python \code{@} réalise le produit matriciel pour
des tableaux à deux dimensions. Dans le script, l'exemple
\code{pytorch\_attention\_example} construit notamment
\code{q = x @ wq}. La matrice \code{wq} y est une matrice
explicite de coefficients ; ne lui attribuons pas automatiquement
l'orientation interne d'un module \code{nn.Linear}.

Chaque ligne d'entrée est transformée par les mêmes colonnes de
coefficients. On peut donc traiter plusieurs positions sans écrire
une boucle Python pour chacune. L'opération ne mélange pas ici
les lignes d'entrée entre elles : chaque ligne produit sa propre
ligne de sortie.

Dans l'exemple chiffré, chaque case demande trois multiplications
et deux additions. Quatre cases nécessitent donc douze multiplications
et huit additions. Ce comptage élémentaire permet déjà d'anticiper
le travail demandé lorsque la longueur des lignes ou le nombre
de sorties augmente.

L'ordre des facteurs importe. Dans notre exemple, \(AR\) a la
forme \(2\times2\), mais \(RA\) aurait la forme \(3\times3\).
Ils ne peuvent même pas être comparés case par case. Avec des
matrices carrées, des formes compatibles ne suffisent toujours
pas à garantir l'égalité des deux produits.

\exercice{Pour \(X=\begin{pmatrix}2&1\\-1&3\end{pmatrix}\) et
\(w=\begin{pmatrix}4\\-2\end{pmatrix}\), calculer \(Xw\).
Quelle opération donnerait plutôt un tableau \(2\times2\) de produits
individuels ?}
\corrige{La première sortie vaut \(2\times4+1\times(-2)=6\).
La seconde vaut \((-1)\times4+3\times(-2)=-10\).
Ainsi \(Xw=(6,-10)^{\mathsf T}\), de forme \(2\times1\).
Un produit élément par élément, éventuellement avec diffusion,
conserverait deux colonnes ; il n'effectuerait pas ces sommes.}

\fiche{8}{Transposée et contrôle des dimensions}{\src{03_attention.py}{MultiHeadSelfAttention.forward}}
\objectif{Permuter des axes volontairement plutôt que corriger les formes au hasard.}

La transposée d'une matrice échange lignes et colonnes. Si \(A\)
a la forme \(m\times n\), alors \(A^{\mathsf T}\) a la forme
\(n\times m\), et sa case \((j,i)\) contient \(A_{ij}\).
Les nombres ne changent pas, seule leur organisation change.
Deux transpositions successives rendent la matrice initiale.
Le nombre total de coefficients est conservé, même si le tableau
devient plus haut ou plus large.
\[
 \begin{pmatrix}1&2&3\\4&5&6\end{pmatrix}^{\mathsf T}
 =\begin{pmatrix}1&4\\2&5\\3&6\end{pmatrix}.
\]
Une ligne devient ainsi une colonne. Cette opération permet de
présenter des vecteurs dans l'orientation nécessaire à un produit
matriciel ; ce n'est ni une inversion ni une division.
Même lorsqu'une matrice inverse existe, elle répond à une autre question.

Supposons que \(Q\) et \(K\) contiennent chacun \(T\) lignes de
\(d_k\) coordonnées. Pour comparer chaque ligne de \(Q\) à chaque
ligne de \(K\), il faut présenter ces dernières comme des colonnes.
Le calcul des formes s'écrit alors :
\[
 \underbrace{Q}_{T\times d_k}
 \underbrace{K^{\mathsf T}}_{d_k\times T}
 =\underbrace{S}_{T\times T},
 \qquad S_{ij}=\sum_{r=0}^{d_k-1}Q_{ir}K_{jr}.
\]
La case \(S_{ij}\) correspond à la ligne \(i\) de \(Q\) et à
la ligne \(j\) de \(K\). En général \(S_{ij}\) n'égale pas
\(S_{ji}\), car \(Q\) et \(K\) sont deux tableaux différents.

\notion{Les axes négatifs se comptent depuis la droite.}
Dans la fonction \code{scaled\_dot\_product\_attention}, appelée
par \code{MultiHeadSelfAttention.forward},
\code{k.transpose(-2, -1)} échange les deux derniers axes.
Les axes précédents, dédiés notamment aux lots et aux têtes,
restent en place. Échanger tous les axes d'un tableau plus grand
serait donc une opération différente.

Une autre règle utile est que la transposée d'un produit inverse
l'ordre des facteurs : \((AR)^{\mathsf T}=R^{\mathsf T}A^{\mathsf T}\).
On peut le vérifier sur chaque case en développant la somme.
Cette inversion d'ordre est indispensable : garder le même ordre
donnerait souvent des dimensions incompatibles, avant même de
considérer les valeurs.

Un contrôle de dimensions doit précéder les valeurs numériques.
Avec \(T=5\) et \(d_k=2\), le produit attendu est
\((5\times2)(2\times5)\), donc \(5\times5\).
Tenter \((5\times2)(5\times2)\) échoue parce que 2 et 5
ne coïncident pas. Si ces nombres étaient accidentellement égaux,
l'absence d'erreur informatique ne prouverait pas la justesse du sens.

\exercice{Soient \(Q=((1,0),(0,2))\) et
\(K=((3,1),(4,-1))\), les couples représentant les lignes.
Calculer \(QK^{\mathsf T}\) et expliquer sa case \((1,0)\).}
\corrige{Les produits scalaires donnent les lignes \((3,4)\)
et \((2,-2)\). La case \((1,0)\) vaut
\(0\times3+2\times1=2\) : elle compare la deuxième ligne
de \(Q\) à la première ligne de \(K\), conformément à
notre convention d'indices commençant à zéro.}

\fiche{9}{Tenseurs : lire les axes B, T et D}{\src{05_model.py}{MiniGPT.forward}}
\objectif{Suivre les tableaux du modèle grâce à la signification de leurs axes.}

Dans ce manuel, un tenseur est un tableau de nombres muni de zéro,
un ou plusieurs axes. Un scalaire n'a aucun axe ; un vecteur en a
un ; une matrice en a deux. Cette définition pratique suffit pour
lire PyTorch, sans introduire la théorie géométrique générale des
tenseurs. Le nombre d'axes et la taille de chaque axe sont deux
informations différentes.

L'entrée \code{x} de \code{MiniGPT.forward} a la forme \([B,T]\).
\(B\) compte les exemples réunis dans le lot ; \(T\) compte les
tokens de chaque exemple. Ainsi \code{x[1, 2]} est l'identifiant du
troisième token du deuxième exemple. Après consultation de la table,
\code{tok\_emb} a la forme \([B,T,D]\), car chaque identifiant
a été remplacé par \(D\) coordonnées.
\[
 [B,T]\longrightarrow[B,T,D]\longrightarrow[B,T,V].
\]
La dernière forme appartient aux logits : chaque position reçoit
un score pour chacun des \(V\) tokens du vocabulaire.
L'axe des caractéristiques \(D\) a été remplacé par l'axe des
classes \(V\), mais les axes des exemples et des positions restent
présents.

\notion{Lire une case avant de compter toutes les cases.}
Dans un tenseur \(h\), \(h_{b,t,j}\) signifie « coordonnée \(j\)
de la position \(t\) dans l'exemple \(b\) ». Cette phrase distingue
les rôles des indices. Pour \(B=2,T=3,D=4\), le tenseur contient
\(2\times3\times4=24\) nombres. Il possède trois axes, pas
vingt-quatre dimensions au sens du nombre d'axes.

La dimension cachée \(D\) ne doit pas être confondue avec le
contexte maximal \(C\). On peut soumettre une séquence plus courte
que \(C\) sans changer \(D\). De même, ajouter des couches \(L\)
augmente le nombre de transformations, pas le nombre de positions.
Le réseau intermédiaire emploie une largeur \(F\), puis revient à \(D\).

Au moment du calcul de la perte, \code{reshape(-1, logits.size(-1))}
réunit les axes \(B,T\) en un axe de \(BT\) prédictions.
Les cibles sont aplaties dans le même ordre. Il faut préserver
cet alignement : avoir autant de lignes que de cibles ne suffit
pas si elles ont été permutées différemment.

\exercice{Avec \(B=3,T=5,D=8,V=11\), donner les formes et
nombres d'éléments des représentations et des logits. Quelle est
la forme des logits aplatis pour la perte ?}
\corrige{Les représentations ont la forme \([3,5,8]\), soit
120 nombres. Les logits ont la forme \([3,5,11]\), soit 165 nombres.
L'aplatissement produit \([15,11]\), accompagné de 15 cibles :
on conserve onze choix pour chacune des quinze prédictions.}

\fiche{10}{Diffusion, view, transpose et contiguous}{\src{03_attention.py}{MultiHeadSelfAttention.forward}}
\objectif{Distinguer une opération sur les valeurs d'une réorganisation des axes et de la mémoire.}

La \emph{diffusion}, ou broadcasting, rend certaines opérations
possibles entre formes différentes. PyTorch aligne les axes depuis
la droite. Sur chaque axe, les tailles doivent être égales, ou l'une
doit valoir un ; les axes absents se comportent comme des axes
de taille un. La valeur singleton est réutilisée aux positions nécessaires.

Dans \code{MiniGPT.forward},
\code{tok\_emb} a la forme \([B,T,D]\) et \code{pos\_emb}
la forme \([1,T,D]\). Leur addition réutilise les mêmes
représentations positionnelles pour chaque exemple. Le résultat
est \([B,T,D]\), sans qu'il soit nécessaire d'écrire \(B\)
copies de la table positionnelle.

\notion{Changer de forme ne signifie pas permuter.}
\code{view} réinterprète une organisation compatible des mêmes
éléments. Il conserve leur nombre et exige des contraintes sur
leur disposition en mémoire. Une suite de 12 valeurs peut devenir
un tableau \([3,4]\), mais pas \([3,5]\). La valeur \(-1\)
demande à PyTorch de déduire une taille ; une seule dimension
peut être ainsi laissée indéterminée.
\[
 [B,T,D]\xrightarrow{\text{view}}[B,T,H,d_k]
 \xrightarrow{\text{transpose}(1,2)}[B,H,T,d_k],
 \qquad D=H d_k.
\]
Ce sont les opérations utilisées pour \code{q}, \code{k} et
\code{v} dans \code{MultiHeadSelfAttention.forward}.
La première sépare les coordonnées en \(H\) groupes ; la seconde
échange les axes des positions et des groupes. Le détail du
calcul des têtes appartient aux fiches sur l'attention.

\code{transpose} peut produire une vue dont les éléments voisins
logiquement ne sont plus voisins dans le stockage. Les pas mémoire,
appelés strides, décrivent alors comment retrouver chaque élément.
Une telle vue n'est pas nécessairement contiguë.

À la sortie, le code utilise
\code{transpose(1, 2).contiguous().view(...)}.
\code{contiguous} garantit une disposition contiguë, en effectuant
une copie si nécessaire, avant la réunion des groupes.
\code{reshape} peut choisir une vue ou une copie ; il ne remplace
cependant jamais une permutation d'axes manquante.

\exercice{Une matrice contient les lignes \((0,1,2)\) et
\((3,4,5)\). Comparer sa transposée et sa réorganisation
contiguë en forme \([3,2]\). Pourquoi le même nombre d'éléments
ne garantit-il pas le même résultat ?}
\corrige{La transposée contient les lignes \((0,3),(1,4),(2,5)\).
Un \code{view(3, 2)} sur la matrice initiale contiguë donne
\((0,1),(2,3),(4,5)\). Les deux résultats contiennent six nombres,
mais ils n'associent pas les mêmes valeurs à chaque indice.
Le premier échange les axes ; le second redécoupe l'ordre existant.}

\fiche{11}{Fonctions affines et orientation de nn.Linear}{\src{04_transformer.py}{FeedForward}}
\objectif{Passer de \(ax+b\) à plusieurs entrées et sorties sans inverser les poids PyTorch.}

Au lycée, une fonction affine s'écrit \(f(x)=ax+b\).
Le coefficient \(a\) règle la variation ; \(b\) fixe la valeur
en zéro. Pour plusieurs entrées, on remplace \(ax\) par une
somme pondérée. Pour plusieurs sorties, on apprend une somme
différente pour chacune.

Un module \code{nn.Linear(D, F)} reçoit \(D\) coordonnées et
produit \(F\) coordonnées. Malgré son nom, il réalise par défaut
une application \emph{affine}, car il contient un biais.
La transformation devient linéaire au sens mathématique si ce biais
est nul ou désactivé.
\[
 y_j=\sum_{i=0}^{D-1}W_{ji}x_i+b_j,
 \qquad W\in\mathbb{R}^{F\times D},\quad b\in\mathbb{R}^{F}.
\]
Attention à l'ordre : \code{weight} est stocké sous la forme
\code{[out\_features, in\_features]}, donc \([F,D]\).
Si les entrées sont les lignes d'une matrice \(X\), le calcul
s'écrit \(Y=XW^{\mathsf T}+b\), avec diffusion du biais sur
les lignes. Écrire \(XW\) utiliserait une autre convention.

\notion{Un exemple à deux sorties.}
Prenons \(x=(2,-1)\), les lignes de poids \((3,4)\) et
\((-2,1)\), puis \(b=(1,-3)\).
La première sortie vaut \(3\times2+4\times(-1)+1=3\).
La seconde vaut \((-2)\times2+1\times(-1)-3=-8\).
Le résultat \((3,-8)\) contient deux scores sans contrainte
de positivité.

Dans \code{FeedForward.net}, la première couche passe de \(D\)
à \(F\), une fonction GELU transforme les valeurs, puis une
seconde couche revient de \(F\) à \(D\). Les mêmes paramètres
s'appliquent à chaque position de chaque exemple. La forme
\([B,T,D]\) devient donc \([B,T,F]\), puis \([B,T,D]\).

Sans transformation non affine entre deux couches affines,
leur composition serait encore affine : développer
\(a_2(a_1x+b_1)+b_2\) donne
\((a_2a_1)x+(a_2b_1+b_2)\).
La fonction intermédiaire apporte donc une capacité que la simple
succession de sommes pondérées n'aurait pas.

Le biais agit indépendamment de l'entrée : si toutes les
coordonnées reçues sont nulles, la sortie est exactement le
vecteur des biais. C'est une manière simple de vérifier son
rôle et son orientation. Une valeur de biais existe par sortie,
pas par entrée ; elle est partagée entre les positions comme
les coefficients de la matrice.

\exercice{Combien de paramètres possède \code{nn.Linear(3, 2)}
avec biais ? Pour une entrée \([4,5,3]\), quelle forme sort
du module ?}
\corrige{La matrice contient \(2\times3=6\) coefficients et
le biais contient deux nombres, soit huit paramètres. La sortie
a la forme \([4,5,2]\). Les vingt positions partagent ces huit
paramètres : leur nombre ne se multiplie pas par le nombre
d'exemples ou de tokens traités.}

\fiche{12}{Exponentielle et stabilité numérique}{\src{03_attention.py}{softmax_manual}}
\objectif{Comprendre pourquoi une formule exacte peut échouer avec des nombres flottants.}

La fonction exponentielle, notée \(\exp(x)=e^x\), associe un
nombre strictement positif à tout réel fini \(x\). Elle vaut un
en zéro, augmente avec \(x\), et transforme les additions en
multiplications : \(e^{a+b}=e^ae^b\). Pour un nombre négatif,
\(e^{-a}=1/e^a\). Ainsi une entrée négative ne produit jamais
une exponentielle négative.

Quelques repères suffisent : \(e\approx2{,}718\),
\(e^2\approx7{,}389\) et \(e^{-2}\approx0{,}1353\).
Une différence additive de un dans l'entrée multiplie la sortie
par environ \(2{,}718\). Cette croissance rapide rend l'exponentielle
utile pour comparer des scores, mais délicate numériquement.

\notion{La machine n'utilise pas tous les réels.}
Un format flottant dispose d'une précision et d'une plage finies.
Une valeur trop grande peut devenir \(+\infty\) : c'est un
débordement. Une valeur positive trop petite peut être arrondie
à zéro. En float32, \(e^{1000}\) déborde largement, bien
que ce nombre existe parfaitement en mathématiques.

La fonction \code{softmax\_manual}, définie dans le script
d'attention, commence par soustraire
le maximum de chaque ligne avant \code{torch.exp}.
Pour les scores \((1000,999)\), elle exponentie donc \((0,-1)\).
\[
 \frac{e^{1000}}{e^{1000}+e^{999}}
 =\frac{1}{1+e^{-1}}\approx0{,}7311.
\]
Le quotient initial pourrait devenir une opération indéfinie
entre infinis dans le programme naïf. Le quotient réécrit utilise
des nombres modestes et donne la même valeur mathématique.
Nous démontrerons l'invariance générale du softmax plus loin.

Après soustraction d'un maximum fini, toutes les exponentielles
sont au plus égales à un et au moins l'une vaut exactement un.
Leur somme ne peut donc pas s'annuler. Des termes très petits
peuvent encore être arrondis à zéro, mais ils ne font plus
déborder le calcul des exponentielles.

La précision limitée affecte aussi les additions : ajouter une
très petite quantité à une très grande peut laisser le résultat
inchangé après arrondi. Réécrire une formule n'est donc pas
seulement une question de présentation.

\attention{La stabilisation suppose une ligne contenant un score
fini et aucun \(+\infty\). Si tous les scores valent
\(-\infty\), soustraire leur maximum produit des valeurs
indéfinies. Masquer toutes les possibilités n'est pas une
distribution de probabilités valide.}

\exercice{Stabiliser les scores \((100,102,101)\), puis donner
les trois exponentielles obtenues et la plus grande probabilité.}
\corrige{On soustrait 102, d'où \((-2,0,-1)\).
Les exponentielles valent environ \((0{,}1353,1,0{,}3679)\).
Leur somme vaut \(1{,}5032\), donc la probabilité maximale
vaut environ \(1/1{,}5032=0{,}6652\), à l'indice 1.
Le décalage n'a pas changé l'ordre des scores.}

\fiche{13}{Logarithme naturel : transformer les produits en sommes}{\src{05_model.py}{MiniGPT.forward}}
\objectif{Donner un sens au logarithme d'une probabilité et à son opposé.}

Le logarithme naturel \(\ln\) est la fonction réciproque de
l'exponentielle. Pour \(x>0\), \(\ln(x)\) est l'unique réel
dont l'exponentielle vaut \(x\). Ainsi \(\ln(1)=0\),
\(\ln(e)=1\) et \(\ln(e^{-2})=-2\).
Son domaine réel exclut zéro et les nombres négatifs.

Cette fonction est croissante : une probabilité plus grande possède
un logarithme plus grand. Pourtant, entre zéro et un, le logarithme
est négatif. Par exemple, \(\ln(0{,}5)\approx-0{,}6931\)
et \(\ln(0{,}1)\approx-2{,}3026\).
Le signe négatif ne signale pas une erreur ; il est la conséquence
normale d'une entrée inférieure à un.
\[
 \ln(ab)=\ln(a)+\ln(b),\qquad
 \ln(a/b)=\ln(a)-\ln(b),\qquad a,b>0.
\]
Ces identités se comprennent en écrivant \(a=e^u\), \(b=e^v\) :
leur produit vaut \(e^{u+v}\), donc son logarithme vaut \(u+v\).
Attention, aucune identité analogue ne donne
\(\ln(a+b)=\ln(a)+\ln(b)\).

\notion{Mesurer la surprise d'une observation.}
Pour une cible observée de probabilité \(p\), la quantité
\(-\ln p\) est positive ou nulle. Elle vaut zéro si \(p=1\)
et augmente fortement quand \(p\) approche zéro.
Cette mesure pénalise donc davantage une cible jugée très improbable.
Elle est exprimée en \emph{nats} lorsqu'on utilise le logarithme
naturel, plutôt qu'en bits.

Dans \code{MiniGPT.forward}, \code{F.cross\_entropy} calcule
une perte liée à ce logarithme, directement à partir des logits.
Le code évite de demander explicitement le logarithme de
probabilités déjà arrondies à zéro. Les formules détaillées de
cette perte apparaîtront après l'étude du softmax.

Pour une suite de prédictions, multiplier beaucoup de probabilités
peut donner un nombre extrêmement petit. Dix probabilités égales
à \(0{,}1\) donnent \(10^{-10}\) ; leurs logarithmes s'additionnent
et donnent \(10\ln(0{,}1)\approx-23{,}026\).
La représentation logarithmique garde ainsi une échelle plus maniable
sans modifier l'ordre de préférence entre probabilités positives.

Le logarithme mesure aussi des améliorations relatives.
Passer d'une probabilité à une probabilité deux fois plus grande
diminue la surprise de \(\ln2\), quelle que soit la valeur
initiale, tant que la nouvelle probabilité reste admissible.
Ce comportement diffère d'une erreur mesurée par simple
soustraction des probabilités : le gain dépend d'un rapport,
pas seulement d'un écart absolu.

\exercice{Deux cibles successives reçoivent les probabilités
\(1/2\) et \(1/4\). Calculer le produit, son logarithme et
la somme des surprises.}
\corrige{Le produit vaut \(1/8\). Comme \(8=2^3\),
\(\ln(1/8)=-3\ln2\approx-2{,}0794\).
La somme des surprises vaut
\(\ln2+\ln4=3\ln2\approx2{,}0794\).
On retrouve exactement l'opposé du logarithme du produit :
les contributions locales s'additionnent.}

\fiche{14}{Probabilités conditionnelles : prévoir avec un contexte}{\src{07_generate.py}{generate}}
\objectif{Interpréter la barre verticale d'une probabilité sans supposer l'indépendance des tokens.}

Une probabilité mesure la plausibilité d'un événement dans une
situation donnée. Pour conditionner sur un événement \(E\) de
probabilité non nulle, on restreint l'univers aux cas où \(E\)
se produit. La probabilité conditionnelle d'un événement \(A\)
s'écrit \(P(A\mid E)\), que l'on lit « probabilité de \(A\)
sachant \(E\) ».
\[
 P(A\mid E)=\frac{P(A\cap E)}{P(E)},
 \qquad P(A\cap E)=P(E)P(A\mid E).
\]
Le symbole \(\cap\) signifie que les deux événements se produisent.
La formule ne dit pas que \(A\) et \(E\) sont indépendants.
L'indépendance serait le cas particulier où connaître \(E\)
ne change pas la probabilité de \(A\).

Imaginons vingt occurrences d'un préfixe dans un petit corpus.
Douze sont suivies du caractère « a », cinq du caractère « e »
et trois d'un espace. Les fréquences conditionnelles observées
sont \(12/20=0{,}6\), \(5/20=0{,}25\) et
\(3/20=0{,}15\). Elles totalisent un parce que ces trois
issues couvrent ici toutes les observations après ce préfixe.

\notion{Le contexte fait partie de la question.}
Le modèle estime une distribution
\(P_\theta(s_{t+1}=v\mid s_0,\ldots,s_t)\), où \(\theta\)
désigne ses paramètres appris et \(v\) un identifiant candidat.
Changer le préfixe change potentiellement toute la distribution.
Les paramètres partagés permettent de construire des estimations
même pour des préfixes jamais vus exactement.

Dans \code{generate}, \code{idx\_cond} garde les derniers tokens
autorisés par \code{context\_length}. Les logits de la dernière
position servent à prévoir le prochain. Si le texte dépasse
la capacité \(C\), ce sont les derniers \(C\) tokens, et non
l'histoire entière, qui sont fournis au modèle sans cache.

Une probabilité prédite n'est pas automatiquement une fréquence
fiable dans le monde réel. Elle dépend du corpus, du modèle
et de l'apprentissage. De même, conditionner sur un texte ne
prouve aucune causalité entre événements du monde : il s'agit
d'une relation de prédiction entre tokens.

Pour chaque contexte fixé, les probabilités de tous les tokens
possibles doivent totaliser un. On ne somme pas les probabilités
d'un même token sur des contextes différents pour effectuer
cette vérification.

\exercice{Dans cent observations, quarante possèdent un préfixe
\(E\), et dix de ces quarante sont suivies du token \(A\).
Calculer \(P(E)\), \(P(A\cap E)\) et \(P(A\mid E)\) à partir
des fréquences. Peut-on en déduire \(P(A)\) ?}
\corrige{On obtient \(P(E)=0{,}4\), \(P(A\cap E)=0{,}1\),
puis \(P(A\mid E)=0{,}1/0{,}4=0{,}25\).
On ne connaît pas les occurrences de \(A\) parmi les soixante
autres observations ; sa probabilité globale reste donc indéterminée.}

\fiche{15}{Loi catégorielle et espérance}{\src{07_generate.py}{generate}}
\objectif{Distinguer un tirage de token, le choix du maximum et une moyenne probabiliste.}

Le prochain token appartient à un ensemble fini de \(V\) catégories.
Une loi catégorielle lui attribue les probabilités
\(p_0,\ldots,p_{V-1}\), toutes positives ou nulles, dont la somme
vaut un. Une réalisation de cette loi est un seul identifiant ;
ce n'est pas le vecteur des probabilités lui-même.

Supposons trois tokens notés « a », « b » et espace, avec les
probabilités \((0{,}5,0{,}3,0{,}2)\). On peut découper
l'intervalle \([0,1[\) en trois morceaux de longueurs correspondantes.
Un nombre uniforme tombant dans le premier morceau choisit « a »,
dans le deuxième « b », et dans le troisième l'espace.
C'est une façon intuitive de comprendre un tirage catégoriel.

Dans \code{generate}, \code{torch.multinomial(probs, num\_samples=1)}
réalise ce type de sélection pour chaque ligne. Malgré le nom
de la fonction, avec un seul tirage par ligne on obtient ici
un échantillon catégoriel. Au contraire, \code{torch.argmax}
choisit systématiquement une catégorie de probabilité maximale.

\notion{L'espérance concerne une quantité numérique associée au résultat.}
Si chaque token \(v\) reçoit une valeur \(g(v)\), sa moyenne
théorique sous la loi est :
\[
 \mathbb{E}[g(S)]=\sum_{v=0}^{V-1}p_vg(v).
\]
Pour \(g(\text{a})=1\), \(g(\text{b})=2\) et
\(g(\text{espace})=0\), elle vaut
\(0{,}5\times1+0{,}3\times2+0{,}2\times0=1{,}1\).
Cette valeur n'est pas un token ; elle résume les résultats
de nombreux tirages.

Faire la moyenne brute des identifiants aurait rarement un sens
linguistique : renuméroter le vocabulaire changerait cette moyenne
sans changer le texte. En revanche, l'indicatrice d'un événement
est utile : elle vaut un si l'événement arrive, zéro sinon.
Son espérance est exactement la probabilité de l'événement.

Sur \(n\) répétitions indépendantes avec la même distribution,
le nombre attendu de « a » vaut \(0{,}5n\). Ce n'est pas
une garantie de résultat exact. Pendant une génération réelle,
la distribution change avec les nouveaux tokens ; on ne peut
donc pas appliquer ce comptage comme si tout le texte provenait
d'une unique loi fixe.

\exercice{Avec les probabilités \((0{,}5,0{,}3,0{,}2)\),
quelle est la probabilité d'obtenir deux fois « b » en deux
tirages indépendants à contexte fixé ? Quel résultat donnerait
le choix du maximum ?}
\corrige{La probabilité vaut \(0{,}3\times0{,}3=0{,}09\).
Le choix du maximum donnerait deux fois « a » sans tirage.
Confondre ces méthodes éliminerait la diversité que le paramètre
\code{sample} permet précisément de conserver.}

\fiche{16}{Chaîne probabiliste autorégressive}{\src{07_generate.py}{generate}}
\objectif{Construire la probabilité d'une continuation à partir de prédictions locales.}

Une génération autorégressive produit les tokens un par un.
Après chaque choix, le token obtenu devient une partie du contexte
du choix suivant. Le préfixe s'allonge donc au cours du calcul.
Cela ne signifie pas que les nouveaux tokens sont indépendants ;
leur dépendance au passé constitue au contraire le principe du modèle.

Pour un préfixe fixé \(c\) et une continuation de \(n\) tokens,
la règle des probabilités conditionnelles donne :
\[
 P_\theta(s_1,\ldots,s_n\mid c)
 =\prod_{t=1}^{n}P_\theta(s_t\mid c,s_1,\ldots,s_{t-1}).
\]
Le signe \(\prod\), lettre grecque pi majuscule, abrège une
multiplication répétée. Il joue pour le produit le rôle du
sigma pour la somme. Au premier facteur, aucun token nouveau
n'a encore été produit : on conditionne seulement sur \(c\).

Supposons que la continuation « a », puis « b » ait les
probabilités successives \(0{,}6\) et \(0{,}25\).
Sa probabilité vaut \(0{,}6\times0{,}25=0{,}15\).
Le deuxième nombre doit être évalué après avoir ajouté « a ».
Utiliser la probabilité de « b » avant cet ajout calculerait
une autre quantité.

\notion{Choix local et meilleure suite sont deux problèmes distincts.}
Imaginons deux premiers choix : « a » reçoit \(0{,}6\),
« b » reçoit \(0{,}4\). Après « a », le meilleur second
choix ne reçoit que \(0{,}5\), tandis qu'après « b », un
second choix reçoit \(0{,}9\). Le meilleur prolongement de
la branche « a » vaut \(0{,}3\), celui de « b » vaut
\(0{,}36\). Choisir le maximum à chaque étape peut donc
manquer la continuation globalement la plus probable.

Dans \code{generate}, \code{torch.cat([idx, next\_token], dim=1)}
réalise l'ajout au préfixe. La boucle recalcule ensuite une
prédiction. Le contexte effectif reste borné par \(C\).
Le produit décrit toujours la loi de la procédure, mais ses
facteurs utilisent alors les suffixes effectivement visibles.

Pour identifier cette loi à celle des logits bruts du modèle,
on considère une température égale à un, sans filtrage.
Modifier la température ou supprimer des candidats modifie les
probabilités utilisées par la procédure de génération.
Ces options seront détaillées après les fondements mathématiques.

\exercice{Trois probabilités conditionnelles d'une continuation
valent \(1/2\), \(1/3\) et \(3/4\). Calculer sa probabilité
et sa surprise totale.}
\corrige{Le produit vaut
\((1/2)(1/3)(3/4)=1/8\). Sa surprise vaut
\(-\ln(1/8)=\ln8=3\ln2\approx2{,}0794\).
Les facteurs dépendent de préfixes successifs différents ;
leur multiplication n'exige aucune hypothèse d'indépendance.}

\fiche{17}{Softmax : définition et invariance par translation}{\src{03_attention.py}{softmax_manual}}
\objectif{Transformer des scores arbitraires en une distribution, puis démontrer une propriété utile.}

Une liste de logits \(z=(z_0,\ldots,z_{V-1})\) peut contenir
des valeurs négatives, positives ou supérieures à un. Pour obtenir
une loi catégorielle, le softmax exponentie chaque score puis divise
par la somme de toutes les exponentielles.
\[
 p_i=\frac{e^{z_i}}{\sum_{j=0}^{V-1}e^{z_j}},
 \qquad 0\leq i<V.
\]
Pour des scores finis, chaque numérateur est strictement positif
et le dénominateur l'est aussi. En additionnant les \(p_i\),
on retrouve le dénominateur divisé par lui-même, donc un.
La fonction est croissante dans l'ordre des scores :
le plus grand logit reçoit la plus grande probabilité.

\notion{Seules les différences entre scores comptent.}
Ajoutons le même réel \(c\) à tous les logits. La propriété
de l'exponentielle permet de factoriser \(e^c\) dans chaque
terme du dénominateur :
\[
 \frac{e^{z_i+c}}{\sum_j e^{z_j+c}}
 =\frac{e^c e^{z_i}}{e^c\sum_j e^{z_j}}
 =p_i.
\]
Le facteur commun disparaît. Les scores \((1,2,3)\) et
\((-2,-1,0)\) définissent donc exactement la même distribution.
Soustraire le maximum, comme le fait \code{softmax\_manual},
est un cas particulier avec \(c=-\max_j z_j\).

Le rapport de deux probabilités vaut
\(p_i/p_j=e^{z_i-z_j}\). Une différence de \(\ln2\)
correspond donc à un rapport de probabilités égal à deux.
Cela fournit une interprétation des écarts entre logits,
mais pas de leur niveau absolu.

Ne confondons pas addition et multiplication : multiplier tous
les logits par deux double leurs écarts et change généralement
la distribution. Cette distinction expliquera l'effet de la
température en génération. De même, augmenter un seul logit
n'est pas une translation commune : sa probabilité augmente
et celles des autres candidats diminuent.

Dans le dépôt, \code{dim=-1} précise l'axe de normalisation.
Pour des logits \([B,T,V]\), le dernier axe est le vocabulaire.
Dans les scores d'attention, il correspond aux positions consultées.
La recette mathématique reste la même, mais les événements
auxquels elle attribue des probabilités ne sont pas les mêmes.

Le dénominateur commun couple les candidats : modifier un score
change également la normalisation des autres. Le softmax n'applique
donc pas une conversion indépendante à chaque case.

\exercice{Les logits \((0,\ln2)\) donnent quelles probabilités ?
Que deviennent-elles après ajout de 7, puis après multiplication
des logits initiaux par deux ?}
\corrige{Les exponentielles initiales valent \((1,2)\), donc
les probabilités sont \((1/3,2/3)\). Ajouter 7 ne change rien.
Multiplier les logits initiaux par deux donne les exponentielles
\((1,4)\), donc \((1/5,4/5)\) : le choix déjà favorisé
devient plus dominant.}

\fiche{18}{Softmax chiffré : du tableau aux probabilités}{\src{07_generate.py}{generate}}
\objectif{Effectuer un softmax exact, contrôler sa somme et interpréter une sélection.}

Prenons un vocabulaire fictif de trois tokens d'identifiants
0, 1 et 2. À la dernière position d'un préfixe, supposons
les logits \(z=(0,\ln2,\ln4)\). Ce choix permet un calcul
exact, car les exponentielles valent des entiers :
\[
 e^z=(1,2,4),\qquad
 \sum_i e^{z_i}=7,\qquad
 p=\left(\frac17,\frac27,\frac47\right).
\]
Les probabilités approchées sont
\((0{,}142857,0{,}285714,0{,}571429)\).
La somme exacte vaut un. Les valeurs décimales arrondies
peuvent présenter un léger écart à un ; il faut distinguer
cet effet d'affichage d'une erreur dans la normalisation.

Refaisons le calcul stabilisé en soustrayant \(\ln4\),
le plus grand logit. Les scores deviennent
\((-\ln4,-\ln2,0)\), leurs exponentielles
\((1/4,1/2,1)\), et leur somme \(7/4\).
Diviser chaque terme par \(7/4\) redonne les mêmes
fractions : \((1/4)/(7/4)=1/7\), puis \(2/7\) et \(4/7\).

\notion{Une probabilité maximale n'est pas une certitude.}
Le token 2 est le candidat le plus probable, mais sa probabilité
reste seulement \(4/7\). Le tirage peut donc retourner 0 ou 1.
Le choix par maximum retourne ici 2 sans hasard.
Le softmax calcule une distribution ; il ne choisit pas lui-même
le token suivant.

Dans \code{generate}, \code{next\_token\_logits} est obtenu
par \code{logits[:, -1, :]}. Si la forme initiale est
\([1,5,3]\), cette sélection donne \([1,3]\) : une seule
distribution sur trois candidats, celle qui suit les cinq tokens
déjà fournis. Les distributions des autres positions ne sont
pas additionnées à celle-ci.

La variable \code{probs} résulte ensuite de \code{torch.softmax}.
Pour retrouver exactement notre exemple, on suppose une température
égale à un et aucun filtrage. Si l'on supprimait le troisième
candidat avant normalisation, les deux probabilités restantes
deviendraient \(1/3\) et \(2/3\), et non \(1/7\) et \(2/7\).
Une distribution doit répartir toute la masse entre les possibilités
encore autorisées.

Pour contrôler un calcul, trois questions suffisent ici :
les valeurs sont-elles non négatives, leur somme vaut-elle un,
et le classement respecte-t-il celui des logits non masqués ?
Ces contrôles ne remplacent pas le calcul exact.

\exercice{Le même calcul est effectué sur les logits
\((\ln3,\ln3,0)\). Donner les probabilités exactes, la somme
des probabilités des deux premiers tokens et le résultat après
soustraction du maximum.}
\corrige{Les exponentielles valent \((3,3,1)\), de somme 7 :
on obtient \((3/7,3/7,1/7)\). Les deux premiers tokens réunis
ont la probabilité \(6/7\). Après soustraction de \(\ln3\),
les exponentielles deviennent \((1,1,1/3)\), de somme \(7/3\).
Le quotient final ne change pas ; les deux premiers candidats
restent ex æquo.}

\fiche{19}{Log-vraisemblance et entropie croisée}{\src{05_model.py}{MiniGPT.forward}}
\objectif{Comprendre la quantité scalaire que l'apprentissage cherche à diminuer.}

Une cible connue est un identifiant \(y\). Si le modèle lui
attribue la probabilité \(p_y\), sa vraisemblance pour cette
observation est \(p_y\). Maximiser cette valeur revient à
maximiser \(\ln p_y\), car le logarithme est croissant.
Pour formuler un problème de minimisation, on prend l'opposé :
la perte vaut \(\ell=-\ln p_y\).
\[
 \ell=-z_y+\ln\!\left(\sum_{j=0}^{V-1}e^{z_j}\right),
 \qquad
 \mathcal{J}=\frac{1}{BT}\sum_{b=0}^{B-1}\sum_{t=0}^{T-1}
 \ell_{b,t}.
\]
La première formule remplace \(p_y\) par son expression softmax.
Elle montre qu'on peut calculer directement la perte à partir
des logits. Le logarithme de la somme d'exponentielles se calcule
de façon stable en retirant leur maximum avant exponentiation.

\notion{Pourquoi parler d'entropie croisée ?}
Représentons la cible par une distribution \(q\) qui vaut un
pour le bon identifiant et zéro ailleurs. L'entropie croisée
est \(-\sum_j q_j\ln p_j\). Tous les termes s'annulent
sauf celui de la cible : on retrouve \(-\ln p_y\).
On n'a pas besoin de fabriquer explicitement ce tableau de zéros
et de un ; les cibles entières suffisent à PyTorch.

Plus généralement, une distribution \(q\) peut répartir la
masse entre plusieurs issues. Son entropie propre est
\(-\sum_j q_j\ln q_j\), tandis que son entropie croisée avec
la prédiction utilise \(\ln p_j\). Les deux expressions
coïncident lorsque \(p=q\). Nous conservons ici les cibles
entières effectivement utilisées dans le dépôt.

Pour \(p=(1/7,2/7,4/7)\), une cible 2 entraîne
\(\ell=\ln(7/4)\approx0{,}5596\). Une cible 0 entraîne
\(\ell=\ln7\approx1{,}9459\). Le même tableau de scores
est donc bon ou mauvais relativement à la réponse observée.

Dans \code{MiniGPT.forward}, \code{F.cross\_entropy} reçoit
les logits de forme \([BT,V]\) et les cibles de forme \([BT]\).
Il ne faut pas appliquer un softmax au préalable : la fonction
attend des scores bruts. La moyenne décrit les positions
évaluées, qui ne sont pas nécessairement des observations indépendantes.

Les cibles proviennent du texte lui-même, grâce au décalage
d'une position. Il n'est donc pas nécessaire qu'une personne
annote chaque prochain caractère. Cette construction explique
le terme d'apprentissage auto-supervisé : le corpus fournit
à la fois les entrées et les réponses attendues.

\exercice{Deux cibles reçoivent les probabilités \(1/2\) et
\(1/8\). Calculer la perte moyenne et la comparer à celle
d'un prédicteur uniforme sur quatre tokens.}
\corrige{La perte moyenne vaut
\((\ln2+\ln8)/2=2\ln2=\ln4\approx1{,}3863\).
Le prédicteur uniforme attribue \(1/4\) à chaque cible et
obtient aussi \(\ln4\). Des erreurs individuelles différentes
peuvent donc conduire à la même perte moyenne.}

\fiche{20}{Dérivée : mesurer une variation locale}{\src{06_train.py}{train}}
\objectif{Retrouver l'idée de pente et comprendre ce que signifie un signal d'amélioration.}

Pour une fonction réelle \(f\), le quotient
\((f(a+h)-f(a))/h\) mesure la variation moyenne entre
deux entrées séparées par \(h\neq0\). Si ce quotient tend
vers une valeur finie lorsque \(h\) tend vers zéro,
cette limite est la dérivée \(f'(a)\).
Elle représente la pente de la tangente au point d'abscisse \(a\).
\[
 f'(a)=\lim_{h\to0}\frac{f(a+h)-f(a)}{h},
 \qquad f(a+h)\approx f(a)+hf'(a).
\]
La deuxième relation est une approximation locale, pas une
égalité pour n'importe quel déplacement. Si \(f'(a)>0\),
un petit déplacement positif augmente la fonction ; un petit
déplacement négatif la diminue. Une dérivée négative inverse
ces indications.

\notion{Un calcul issu de la définition.}
Pour \(f(x)=x^2\), développer \((a+h)^2\) donne
\(a^2+2ah+h^2\). Le quotient de variation vaut donc
\(2a+h\), et sa limite vaut \(2a\). Au point \(a=3\),
la dérivée vaut 6. Un déplacement \(h=0{,}01\) prédit
une hausse d'environ \(0{,}06\) ; la hausse exacte vaut
\(3{,}01^2-9=0{,}0601\).

Les règles utiles pour la suite sont :
la dérivée d'une constante est zéro, celle de \(ax+b\)
est \(a\), celle de \(e^x\) est \(e^x\), et celle
de \(\ln x\), pour \(x>0\), est \(1/x\).
La dérivée d'une somme est la somme des dérivées.
Ces quelques règles couvrent une grande partie des opérations
de nos petits exemples.

Dans \code{train}, \code{loss.backward()} calcule des
dérivées de la perte par rapport aux paramètres.
Il ne remplace pas chaque poids par une infinité de petites
perturbations : il utilise les règles de dérivation des opérations
composées. Nous verrons ensuite comment cela se fait.

Une pente nulle ne prouve pas qu'un point est un minimum.
Pour \(f(x)=-x^2\), la dérivée en zéro est nulle,
mais zéro est un maximum. L'information de premier ordre
est locale et doit être interprétée avec prudence.

On peut estimer une dérivée en comparant deux valeurs proches,
mais choisir un écart minuscule n'est pas toujours plus fiable
sur ordinateur. Les arrondis peuvent masquer la différence entre
les deux résultats. La dérivation des opérations évite cette
soustraction fragile et fournit directement les sensibilités
des calculs effectivement composés.

\exercice{Pour \(f(w)=(w-3)^2\), calculer \(f'(1)\).
Prévoir l'effet d'un déplacement de \(1\) vers \(1{,}1\),
puis comparer avec les valeurs exactes.}
\corrige{On développe \(f(w)=w^2-6w+9\), donc
\(f'(w)=2w-6\) et \(f'(1)=-4\).
La variation prévue est \(0{,}1\times(-4)=-0{,}4\).
La perte passe exactement de 4 à \(3{,}61\), soit
une variation de \(-0{,}39\) : l'approximation indique
correctement le sens de l'amélioration.}

\fiche{21}{Règle de chaîne : dériver une composition}{\src{04_transformer.py}{FeedForward.forward}}
\objectif{Propager une petite variation à travers plusieurs calculs successifs.}

Une composition applique une fonction au résultat d'une autre.
Si \(u=g(x)\) puis \(y=f(u)\), on écrit
\(y=f(g(x))\). Une petite variation de \(x\) modifie
d'abord \(u\), puis cette modification se répercute sur \(y\).
La règle de chaîne multiplie ces deux sensibilités locales :
\[
 \frac{dy}{dx}=\frac{dy}{du}\frac{du}{dx},
 \qquad (f\circ g)'(x)=f'(g(x))g'(x).
\]
La notation ressemble à une simplification de fractions,
mais elle exprime une règle de dérivation. Le point auquel
chaque dérivée est évaluée compte : \(f'\) s'évalue à
la valeur intermédiaire \(g(x)\), pas automatiquement à \(x\).
Cette règle suppose que les fonctions rencontrées soient dérivables
aux points utilisés pendant le calcul.

Considérons \(u=3x-1\), puis \(y=u^2\).
On a \(du/dx=3\) et \(dy/du=2u\), donc
\(dy/dx=6(3x-1)\). À \(x=2\), l'intermédiaire
vaut 5 et la dérivée finale vaut 30.
Développer \(y=9x^2-6x+1\) donne également
\(18x-6=30\), ce qui vérifie le calcul par une autre voie.

\notion{Les variables intermédiaires rendent les calculs lisibles.}
Au lieu de dériver une expression immense, on conserve les
résultats de chaque opération. Pour
\(a=wx+b\), \(h=\phi(a)\), \(z=vh+c\), la sensibilité
de \(z\) à \(w\), les autres valeurs étant fixées, est
\(v\phi'(a)x\). On retrouve les facteurs du trajet
\(w\rightarrow a\rightarrow h\rightarrow z\).

\code{FeedForward.forward} applique une séquence :
couche affine, GELU, couche affine, puis dropout.
Les fonctions réellement utilisées ont leurs règles de dérivation.
Notre carré \(u^2\) est seulement un exemple facile à calculer,
pas la définition de GELU.

Si plusieurs trajets relient une variable à la sortie,
leurs contributions s'ajoutent. Pour \(y=x^2+x\),
la branche carrée apporte \(2x\) et la branche directe
apporte 1. La dérivée totale est \(2x+1\).
Cette idée éclairera plus tard les additions résiduelles
du Transformer.

La variable par rapport à laquelle on dérive doit toujours
être annoncée. Dans \(a=wx+b\), dériver par rapport à
\(x\) donne \(w\), tandis que dériver par rapport à
\(w\) donne \(x\). Le biais est constant dans ces deux
calculs, mais sa dérivée vaut un lorsqu'on dérive par rapport
à \(b\). Une même expression possède donc plusieurs sensibilités,
selon la question posée.

Des produits répétés de petites sensibilités peuvent atténuer
fortement une variation ; de grandes sensibilités peuvent
l'amplifier. Ce constat arithmétique prépare l'étude des
réseaux profonds, sans remplacer l'analyse de leurs valeurs réelles.

\exercice{Dériver \(f(x)=\ln(1+e^x)\) en nommant ses
intermédiaires. Quelle est la dérivée en zéro ?}
\corrige{Posons \(u=e^x\), \(v=1+u\), puis \(f=\ln v\).
Les sensibilités sont \(du/dx=e^x\), \(dv/du=1\) et
\(df/dv=1/v\). Leur produit vaut
\(f'(x)=e^x/(1+e^x)\). En zéro, \(e^0=1\),
donc \(f'(0)=1/2\).}

\fiche{22}{Gradient : une dérivée pour chaque paramètre}{\src{06_train.py}{train}}
\objectif{Étendre la pente d'une courbe à une perte dépendant de plusieurs nombres.}

Un réseau possède de nombreux paramètres, regroupés symboliquement
dans \(\theta=(\theta_0,\ldots,\theta_{n-1})\).
Sa perte \(\mathcal{J}(\theta)\) reste un seul nombre.
La dérivée partielle
\(\partial\mathcal{J}/\partial\theta_i\) mesure l'effet local
du paramètre \(i\) lorsque les autres sont maintenus fixes.
Le symbole \(\partial\) rappelle cette convention.

Le gradient rassemble toutes ces dérivées dans le même ordre
que les paramètres. Il indique comment une petite modification
simultanée se répercute au premier ordre :
\[
 \nabla\mathcal{J}=
 \left(\frac{\partial\mathcal{J}}{\partial\theta_0},
 \ldots,\frac{\partial\mathcal{J}}{\partial\theta_{n-1}}\right),
 \qquad
 \Delta\mathcal{J}\approx\nabla\mathcal{J}\cdot\Delta\theta.
\]
Nous retrouvons le produit scalaire : chaque variation est
pondérée par sa sensibilité, puis toutes les contributions
s'additionnent.

Pour \(\mathcal{J}(a,b)=(a+2b-3)^2\), la règle de chaîne
donne les dérivées \(2(a+2b-3)\) et \(4(a+2b-3)\).
Au point \((0,0)\), le gradient est \((-6,-12)\).
Une perturbation \((0{,}01,0)\) prévoit une baisse de
\(0{,}06\) ; la perturbation \((0,0{,}01)\) prévoit
une baisse de \(0{,}12\). Les deux paramètres n'ont pas
la même influence locale.

\notion{Le gradient des logits possède une forme particulièrement simple.}
Pour la perte d'une cible \(y\) étudiée précédemment,
la dérivation du logarithme de la somme donne :
\[
 \frac{\partial\ell}{\partial z_i}
 =p_i-\mathbf{1}_{i=y}.
\]
L'indicatrice \(\mathbf{1}_{i=y}\) vaut un pour la cible,
zéro sinon. Le terme \(p_i\) vient de
\(e^{z_i}/\sum_j e^{z_j}\), et la dérivée de \(-z_y\)
fournit le terme négatif. Augmenter le logit correct diminue
ainsi localement la perte, tandis qu'augmenter un concurrent
l'augmente.

Dans \code{train}, \code{loss.backward()} propage ces
sensibilités jusqu'aux paramètres. Leur attribut \code{grad}
possède leur forme : une matrice de poids reçoit une matrice
de dérivées. Pour une perte moyenne sur plusieurs positions,
les contributions sont additionnées puis divisées par leur nombre.
Les identifiants entiers ne sont pas les paramètres à ajuster.

Un poids partagé intervient à plusieurs positions. Sa dérivée
totale additionne les contributions de toutes ces utilisations :
on n'apprend pas une copie indépendante du poids pour chaque
token. Cette accumulation explique aussi pourquoi les anciennes
valeurs de \code{grad} doivent être effacées lorsqu'on commence
une nouvelle mise à jour sans accumulation volontaire.

\exercice{Avec \(p=(0{,}2,0{,}5,0{,}3)\) et une cible
d'identifiant 2, calculer le gradient par rapport aux logits,
puis la somme de ses coordonnées.}
\corrige{Le gradient vaut \((0{,}2,0{,}5,-0{,}7)\),
de somme zéro. Une translation commune de tous les logits
a donc une variation de perte nulle au premier ordre.
Cela concorde avec l'invariance exacte du softmax démontrée
précédemment.}

\fiche{23}{Rétropropagation : un petit réseau calculé entièrement}{\src{06_train.py}{train}}
\objectif{Effectuer un passage avant et un passage arrière sans traiter la dérivation automatique comme une magie.}

Construisons un réseau scalaire pédagogique. Son entrée est \(x\),
ses paramètres sont \(w,b,v,c\), et ses deux logits sont
\((z,0)\). La cible est le premier token, d'identifiant zéro.
La fonction carrée remplace volontairement GELU pour permettre
un calcul intégral à la main ; l'entropie croisée reste celle
étudiée pour le modèle.
\[
 a=wx+b,\qquad h=a^2,\qquad z=vh+c,\qquad
 \ell=\ln(1+e^{-z}).
\]
La dernière formule vient de
\(-\ln(e^z/(e^z+1))\). Le logit du deuxième token reste
fixé à zéro : ce n'est pas un paramètre supplémentaire.
Nous choisissons \(x=2,w=1,b=-1,v=2,c=-2\).

\notion{Passage avant : mémoriser les intermédiaires.}
On calcule successivement \(a=1\times2-1=1\),
\(h=1^2=1\), puis \(z=2\times1-2=0\).
Les deux logits sont donc nuls, les probabilités valent
\((1/2,1/2)\) et la perte vaut \(\ln2\approx0{,}6931\).
Ces valeurs intermédiaires serviront dans les dérivées.

Le passage arrière commence à la perte.
La dérivée par rapport à \(z\) vaut
\(-1/(1+e^z)\), donc \(-1/2\) ici.
Pour l'opération \(z=vh+c\), on multiplie ce signal par
la dérivée locale correspondant à chaque entrée :
\(\partial\ell/\partial v=(-1/2)h=-1/2\),
\(\partial\ell/\partial c=-1/2\), et
\(\partial\ell/\partial h=(-1/2)v=-1\).

À travers le carré, on obtient
\(\partial\ell/\partial a=(-1)\times2a=-2\).
La première couche affine reçoit alors ce signal :
\[
 \frac{\partial\ell}{\partial w}=-2x=-4,\qquad
 \frac{\partial\ell}{\partial b}=-2,\qquad
 \frac{\partial\ell}{\partial x}=-2w=-2.
\]
Le gradient des quatre paramètres, dans l'ordre \((w,b,v,c)\),
est donc \((-4,-2,-1/2,-1/2)\).
La dérivée de l'entrée a été calculée pour suivre toute la chaîne,
mais l'entrée observée ne doit pas être modifiée pour apprendre
les paramètres.

Dans \code{train}, \code{loss.backward()} automatise ce
parcours inverse pour le vrai graphe du MiniGPT.
Chaque opération transmet une sensibilité locale ; les chemins
partagés additionnent leurs contributions. Le passage arrière
ne change pas encore les paramètres : cette responsabilité
appartient à l'optimiseur.

La valeur initiale du signal arrière est un :
la dérivée de la perte par rapport à elle-même vaut un.
Chaque étape conserve ensuite les paramètres aux valeurs utilisées
pendant le passage avant. Si l'on changeait un poids au milieu
de ce parcours, les dérivées locales ne décriraient plus toutes
le même réseau et le gradient obtenu perdrait sa cohérence.

\exercice{En ne changeant que \(c\), appliquer un déplacement
opposé au gradient de taille \(0{,}1\) fois sa valeur.
Recalculer le logit et la perte.}
\corrige{Le nouveau biais vaut
\(-2-0{,}1(-1/2)=-1{,}95\). Le logit devient
\(z=2-1{,}95=0{,}05\). La perte vaut
\(\ln(1+e^{-0{,}05})\approx0{,}6685\), inférieure à
\(\ln2\), conformément au signe du gradient.}

\fiche{24}{Descente de gradient et synthèse des fondements}{\src{06_train.py}{train}}
\objectif{Relier la prédiction, la perte et les dérivées à une règle élémentaire d'apprentissage.}

Une fois le gradient connu, comment modifier les paramètres ?
La descente de gradient élémentaire avance dans la direction
opposée au gradient. Le taux d'apprentissage \(\eta>0\)
règle la taille du déplacement :
\[
 \theta_{\mathrm{nouveau}}
 =\theta_{\mathrm{ancien}}-\eta\nabla\mathcal{J}
   (\theta_{\mathrm{ancien}}).
\]
Tous les composants utilisent les dérivées évaluées aux anciennes
valeurs. Les remplacer successivement en recalculant une partie
du gradient définirait une autre procédure.
Un paramètre dont la dérivée est négative augmente ; soustraire
un nombre négatif revient à l'ajouter.

\notion{Pourquoi cette direction peut-elle améliorer la perte ?}
L'approximation locale donne, pour le déplacement
\(\Delta\theta=-\eta\nabla\mathcal{J}\),
une variation d'environ
\(-\eta\|\nabla\mathcal{J}\|^2\).
Elle est négative si le gradient est non nul.
Le raisonnement reste local : un pas trop grand peut sortir
de la région où cette approximation est pertinente.

Sur \(\mathcal{J}(w)=(w-3)^2\), partons de \(w=1\).
Le gradient vaut \(-4\). Avec \(\eta=0{,}1\), on obtient
\(w_{\mathrm{nouveau}}=1{,}4\) et une perte \(2{,}56\),
contre 4 auparavant. Avec \(\eta=1{,}5\), on obtient 7
et une perte 16 : suivre la bonne direction initiale
n'empêche pas de dépasser une région favorable.

Dans \code{train}, le programme emploie
\code{torch.optim.AdamW}, non cette descente élémentaire.
AdamW adapte ses mises à jour à l'historique des gradients
et applique une décroissance des poids. La formule précédente
explique l'intuition, mais ne reproduit pas ses étapes exactes.
Les détails de l'entraînement seront développés dans les
fiches 73 à 96.

L'ordre effectif du code est important :
le passage avant calcule les logits et la perte,
\code{optimizer.zero\_grad(set\_to\_none=True)} efface
les anciennes dérivées, \code{loss.backward()} calcule
les nouvelles, puis \code{optimizer.step()} modifie
les paramètres. Le calcul des gradients et leur application
constituent deux opérations distinctes.

Nous disposons désormais de la chaîne mathématique essentielle :
les indices désignent les données ; les vecteurs et matrices
les représentent ; les fonctions composées donnent des logits ;
le softmax fournit des probabilités ; l'entropie croisée
mesure leur accord avec les cibles ; les dérivées indiquent
comment modifier les paramètres. Les fiches suivantes pourront
donc préciser les données et les représentations avant
l'attention, réservée aux fiches 49 à 72.

\exercice{Une perte a pour gradient \((2,-3)\) au point
\((1,1)\). Avec \(\eta=0{,}05\), calculer le nouveau point
et la variation de perte prédite au premier ordre.}
\corrige{Le déplacement vaut \((-0{,}1,0{,}15)\), donc
le nouveau point est \((0{,}9,1{,}15)\).
La variation prévue est
\(2(-0{,}1)+(-3)(0{,}15)=-0{,}65\).
C'est une prévision locale, pas une garantie que la perte
réelle diminue exactement de cette quantité.}
